\section{Results}
\subsection{Controlled Visual Test and Supervisory-Loop Timing}
The separate controlled closed-set visual test used 180 photographs, with 30 images per object class, acquired under fixed viewpoint, background, and lighting conditions. All 180 images were classified correctly and mapped to the intended manipulation strategy. The mean confidence score was 0.853, and the mean per-image processing time was 48.19 ms (Fig.~\ref{fig:4}). Class-wise results and the confusion matrix are provided in Supplementary Table~S3 and Supplementary Fig.~S1.

These results characterize closed-set recognition and strategy mapping for the six tested object instances under the specified controlled imaging conditions. They do not quantify the correctness of online strategy triggering, trigger-to-contact timing, or performance on untested objects.

\begin{figure*}[!t]
\centering
\includegraphics[width=\textwidth]{fig4.png}
\caption{Confidence scores and per-image processing time in the controlled closed-set visual test ($n=30$ images per object class). Bars indicate mean $\pm$ SD; half violins, points, segments, and open circles indicate distributions, individual observations, IQRs, and medians, respectively. Reference lines indicate the confidence threshold of 0.25 and the overall means for confidence (0.853) and processing time (48.19 ms). The reported processing time does not represent end-to-end task-trigger latency.}
\label{fig:4}
\end{figure*}

The median logged duration of the nominal 200-Hz supervisory cycle was 5.07 ms.

\subsection{Overall Results Across the Five Modes}
Table~\ref{tab:8} provides descriptive context across all five modes; the primary paired C--E evidence is reported separately in Table~\ref{tab:9} and Fig.~\ref{fig:7}.
\begin{table*}[!t]
\centering
\caption{Descriptive task-performance and workload results across the five experimental modes.}
\label{tab:8}
\small
\setlength{\tabcolsep}{5pt}
\renewcommand{\arraystretch}{1.16}
\begin{tabularx}{\textwidth}{@{}p{0.11\textwidth}YYYp{0.15\textwidth}@{}}
\toprule
Mode & Task completion time (s) & Master-side trajectory length (m) & Observed success, $n/N$ (\%) & Raw NASA-TLX \\
\midrule
A & 21.18 [20.62, 22.08] & 0.757 [0.693, 0.816] & 22/27 (81.5\%) & 62.50 [59.67, 64.50] \\
B & 20.89 [20.12, 21.83] & 0.787 [0.721, 0.861] & 21/27 (77.8\%) & 56.17 [55.00, 59.33] \\
\textbf{C} & \textbf{19.57 [18.41, 20.05]} & \textbf{0.697 [0.660, 0.769]} & \textbf{26/27 (96.3\%)} & \textbf{48.67 [47.67, 51.83]} \\
D & 20.79 [20.32, 21.16] & 0.722 [0.678, 0.768] & 24/27 (88.9\%) & 60.33 [57.33, 62.50] \\
E & 20.73 [19.95, 22.25] & 0.732 [0.678, 0.799] & 24/27 (88.9\%) & 53.67 [51.83, 57.83] \\
\bottomrule
\end{tabularx}
\par\smallskip
\footnotesize\textit{Note:} Continuous outcomes are reported as median [IQR], whereas observed success is reported as successful trials/total trials (\%). Task completion time, master-side trajectory length, and observed success use $n=27$ trial-level observations per mode. Raw NASA-TLX uses $n=9$ operator $\times$ strategy questionnaire units per mode. Mode definitions are provided in Table~\ref{tab:3}, and the trial-level task-performance data are provided in Supplementary Table~S4.
\end{table*}
Across the five modes, mode C showed the lowest median completion time, the highest observed success rate, the lowest median Raw NASA-TLX score, and the lowest median master-side trajectory length (Table~\ref{tab:8} and Figs.~\ref{fig:5}--\ref{fig:6}). Using the trial-level data in Supplementary Table~S4, the mean completion time was 19.28 $\pm$ 1.30 s in mode C, compared with 21.42 $\pm$ 1.58 s in mode A, 21.01 $\pm$ 1.61 s in mode B, 20.91 $\pm$ 1.10 s in mode D, and 21.07 $\pm$ 1.56 s in mode E. With each comparison mode used as the reference, the corresponding mean reductions for mode C were 10.0\%, 8.2\%, 7.8\%, and 8.5\%, respectively. Consistent with the predefined outcome hierarchy, inferential analyses focused on task completion time as the primary endpoint.

Using the 27 operator-nested matched task blocks as repeated-measures units, the Friedman test indicated an overall difference in completion time across the five modes ($\chi^2(4)=30.904$, $p<0.001$). Holm-adjusted Wilcoxon signed-rank tests indicated lower completion times in mode C than in modes A, B, D, and E, with all adjusted $p$-values below 0.01. These comparisons were interpreted as conditional repeated-measures evidence for the observed three-operator sample. Descriptively, mode C yielded the lowest Raw NASA-TLX score within each of the nine operator $\times$ strategy questionnaire units.

\begin{figure}[!t]
\centering
\includegraphics[width=\columnwidth]{fig5.png}
\caption{Task completion time across the five experimental modes ($n=27$ matched task blocks per mode). Filled circles, thick horizontal segments, whiskers, and open circles indicate medians, IQRs, 1.5$\times$IQR ranges, and outliers, respectively. This descriptive panel displays marginal mode distributions; the paired C--E observations are shown in Fig.~\ref{fig:7}.}
\label{fig:5}
\end{figure}

\begin{figure*}[!t]
\centering
\includegraphics[width=\textwidth]{fig6.png}
\caption{Descriptive workload and observed task-success results across the five experimental modes. (a) Raw NASA-TLX scores: connected markers show each operator's mean across the three strategy-level questionnaire units within each mode ($n=9$ operator $\times$ strategy units per mode), and marker shape identifies the operator. (b) Number of successful trials out of 27 trials per mode. Both panels are descriptive; no confidence intervals or primary inferential tests are shown.}
\label{fig:6}
\end{figure*}

\subsection{Primary C--E Comparison and Exploratory Process Measures}
The primary paired comparison between modes C and E is summarized in Table~\ref{tab:9} and Fig.~\ref{fig:7}.

\begin{table*}[!t]
\centering
\caption{Primary within-sample comparison of complete multi-channel configuration and impedance-only reconfiguration.}
\label{tab:9}
\small
\setlength{\tabcolsep}{5pt}
\renewcommand{\arraystretch}{1.16}
\begin{tabularx}{\textwidth}{@{}p{0.18\textwidth}YYYp{0.20\textwidth}@{}}
\toprule
Outcome & Mode C, median [IQR] & Mode E, median [IQR] & Mean paired difference $\Delta$ (E--C) [95\% CI] & Operator-level direction \\
\midrule
Task completion time (s) & 19.57 [18.41, 20.05] & 20.73 [19.95, 22.25] & 1.79 [1.10, 2.51] & 3/3 favored C \\
Master-side trajectory length (m) & 0.697 [0.660, 0.769] & 0.732 [0.678, 0.799] & 0.024 [-0.014, 0.059] & Mixed \\
Raw NASA-TLX & 48.67 [47.67, 51.83] & 53.67 [51.83, 57.83] & 4.87 [---] & 3/3 favored C \\
\bottomrule
\end{tabularx}
\par\smallskip
\footnotesize\textit{Note:} Differences were defined as mode E minus mode C, so positive $\Delta$ values favor mode C for all three outcomes. Task completion time and master-side trajectory length use $n=27$ matched task blocks. Raw NASA-TLX uses $n=9$ paired operator $\times$ strategy questionnaire units per mode. Confidence intervals for completion time and trajectory length are matched-block bootstrap 95\% intervals based on 10{,}000 resamples and are conditional on the observed block structure. No bootstrap confidence interval was calculated for Raw NASA-TLX.
\end{table*}
Median task completion time was 19.57 s in mode C and 20.73 s in mode E. Defining the block-level difference as $\Delta T=T_E-T_C$, the mean paired difference was 1.79 s, with a matched-block bootstrap 95\% confidence interval of 1.10--2.51 s. Relative to the mean completion time in mode E, this difference corresponded to an 8.5\% reduction in mode C. Mode C had a lower completion time than mode E in 22 of the 27 matched task blocks.

A successful-trial sensitivity analysis (23 blocks with both trials successful) yielded a similar estimate (mean paired difference 1.80 s; 19/23 blocks favoring mode C).

For the paired success outcomes, 23 matched blocks were successful in both modes, three were successful in mode C only, one was successful in mode E only, and none failed in both modes. The exact two-sided McNemar test yielded $p=0.625$ and, given only four discordant pairs, did not provide clear evidence of a mode-related difference in observed success. Accordingly, the observed success-rate difference was treated as supplementary descriptive evidence.

Median master-side trajectory length was 0.697 m in mode C and 0.732 m in mode E. With $\Delta L=L_E-L_C$, the mean paired trajectory-length difference was 0.024 m, with a matched-block bootstrap 95\% confidence interval of -0.014 to 0.059 m. Because the interval included zero, the data did not provide clear evidence of a C--E difference in master-side trajectory length. Pause counts were lower on average in mode C, although the two modes had the same median pause count.

\begin{figure}[!t]
\centering
\includegraphics[width=\columnwidth]{fig7.png}
\caption{Paired task completion times for modes C and E across 27 matched task blocks. Points below the identity line indicate shorter completion time in mode C. For the paired difference $\Delta T=T_E-T_C$, positive values favor mode C. The violin, horizontal segment, and diamond show the distribution, median, and mean of the paired differences, respectively. The mean paired difference was 1.79 s (matched-block bootstrap 95\% CI, 1.10--2.51 s), and 22 of 27 blocks favored mode C.}
\label{fig:7}
\end{figure}

Using the predefined pause criterion in Section~3.3, mode C had a median [IQR] pause count of 3 [2, 3.5] and a mean of 2.74 $\pm$ 1.23, whereas mode E had a median pause count of 3 and a mean of 3.41 $\pm$ 1.67. The mean pause count was lower in mode C within each of the three manipulation strategies, while the overall median remained 3 in both modes. Pause count was treated descriptively, and no formal hypothesis test was performed. The robustness of the result to alternative speed or duration thresholds was not evaluated.

\subsection{Observed Failure Outcomes}
Across the 135 trials, 18 were classified as failures: 5, 6, 1, 3, and 3 in modes A--E, respectively. The observed failure outcomes involved object drops, observable slip, or visible damage according to the predefined criteria in Section~3.3. Table~\ref{tab:10} summarizes these descriptive mode-level counts and representative trial observations.

\begin{table*}[!t]
\centering
\caption{Observed failure counts and representative trial observations by experimental mode.}
\label{tab:10}
\small
\setlength{\tabcolsep}{6pt}
\renewcommand{\arraystretch}{1.16}
\begin{tabularx}{\textwidth}{@{}p{0.19\textwidth}p{0.14\textwidth}Y@{}}
\toprule
Mode & Observed failures, $n/N$ & Representative trial observations \\
\midrule
A & 5/27 & Paper-cup deformation and difficulty maintaining the scissors' pose were noted in failed trials. \\
B & 6/27 & No consistent failure pattern could be assigned from the retained trial records. \\
\textbf{C} & \textbf{1/27} & Observable mouse slip during transfer. \\
D & 3/27 & Object drops or loss of grasp after grasp establishment. \\
E & 3/27 & Observed drops, slip, or visible damage during grasp establishment or transfer. \\
\bottomrule
\end{tabularx}
\par\smallskip
\footnotesize\textit{Note:} Mode definitions are provided in Table~\ref{tab:3}. Failure labels were assigned during the experiment using the predefined criteria in Section~3.3. The listed observations are representative rather than exhaustive. Systematic video recording and independent blinded re-evaluation were not performed; the observations should not be interpreted as confirmed root-cause classifications.
\end{table*}
Among the five modes, mode C had the lowest observed failure count, with 1 failure in 27 trials.

\subsection{Directional Consistency Across Operators and Objects}
The mean paired C--E completion-time difference favored mode C for each of the three operators and each of the six tested objects.

At the operator level, the mean paired differences $\Delta T=T_E-T_C$ were 1.66, 2.56, and 1.16 s for P01, P02, and P03, respectively (mode C faster in 7/9, 9/9, and 6/9 matched blocks). Leave-one-operator-out mean paired differences were 1.86, 1.41, and 2.11 s, all remaining positive. Operator-level medians [IQR] are shown in Fig.~\ref{fig:8}.

At the object level, mode C had the lowest mean task completion time among the five modes for each of the six tested objects (Supplementary Table~S5). Relative to mode E, the mean completion-time reduction in mode C ranged from 3.3\% for the bottle to 13.2\% for the scissors. The mean C--E difference favored mode C for all six objects, while the magnitude varied across objects.
